% Compile with: uplatex file.tex && dvipdfmx file.dvi
%           or: latexmk -uplatex file.tex
\documentclass[dvipdfmx]{article}
\usepackage{amsmath, amssymb}
\usepackage{tikz}

\title{scenario\_information\_theory}
\date{\today}

\begin{document}
\maketitle

Write the Shannon entropy $H=-\sum_x p(x)\log_2 p(x)$ for a fair die ($p=1/6$, six faces) as the definition's own sum, in bits. No dedicated operation is used.

$H = -\sum_{x=1}^{6} \frac{1}{6}\log_2 \frac{1}{6}$

\begin{equation}
  H = -\sum_{x=1}^{6} \frac{1}{6}\log_2 \frac{1}{6}
  \label{eq:E1}
\end{equation}

The right-hand side evaluates to $\log 6/\log 2 = \log_2 6 \approx 2.585$ bits — information theory's default unit, so this is the form the definition should be written in. Until 2026-08-16 the based spelling was silent-wrong: it became `log(2)*6`, the logarithm of 2 times 6.

compute the right-hand side of E1

\begin{equation}
  -\sum_{x=1}^{6} \frac{1}{6}\log_2 \frac{1}{6} = \frac{\log 6}{\log 2}
  \label{eq:E2}
\end{equation}

Next the entropy of a Bernoulli distribution, $H_b(p)$. From here the general operations compose naturally.

$H_b = -p\log p - (1-p)\log(1-p)$

\begin{equation}
  H_b = -p\log p - (1-p)\log(1-p)
  \label{eq:E3}
\end{equation}

Differentiate with respect to $p$ to look for the condition that maximises the entropy.

differentiate E3 with respect to p

\begin{equation}
  \frac{dH_b}{dp} = \log \left(1-p\right)-\log p
  \label{eq:E4}
\end{equation}

Set the derivative to zero and solve for $p$. The CAS returns $\log p=\log(1-p)$ — it does not take the injectivity-of-log step, so $p=1/2$ is read off by the researcher. Not filling that in is the correct behaviour.

solve E4 for p

\begin{equation}
  \log p = \log \left(1-p\right)
  \label{eq:E5}
\end{equation}

Write the Kullback–Leibler divergence between two Bernoulli distributions straight from the definition. The point is that the summand's $\log\frac{p}{q}$ can simply be typed.

$D = p\log \frac{p}{q} + (1-p)\log \frac{1-p}{1-q}$

\begin{equation}
  D = p\log \frac{p}{q} + (1-p)\log \frac{1-p}{1-q}
  \label{eq:E6}
\end{equation}

Substitute $p$ for $q$. The basic KL property $D(p\Vert p)=0$ comes out as a computation, not as a displayed definition.

substitute p for q in E6

\begin{equation}
  D = 0
  \label{eq:E7}
\end{equation}

Write the mutual information as the definition's own double sum.

$I = \sum_{x=1}^{2}\sum_{y=1}^{2} p_{xy}(x,y)\log \frac{p_{xy}(x,y)}{p_x(x) p_y(y)}$

\begin{equation}
  I = \sum_{x=1}^{2}\sum_{y=1}^{2} p_{xy}(x,y)\log \frac{p_{xy}(x,y)}{p_x(x) p_y(y)}
  \label{eq:E8}
\end{equation}

Write the independent case, i.e. with $p_{xy}(x,y)=p_x(x)p_y(y)$ substituted in.

$I_0 = \sum_{x=1}^{2}\sum_{y=1}^{2} p_x(x) p_y(y)\log \frac{p_x(x) p_y(y)}{p_x(x) p_y(y)}$

\begin{equation}
  I_0 = \sum_{x=1}^{2}\sum_{y=1}^{2} p_x(x) p_y(y)\log \frac{p_x(x) p_y(y)}{p_x(x) p_y(y)}
  \label{eq:E9}
\end{equation}

Evaluating the right-hand side gives 0: independence makes the mutual information vanish, verified by computation.

compute the right-hand side of E9

\begin{equation}
  \sum_{x=1}^{2}\sum_{y=1}^{2} p_x(x) p_y(y)\log \frac{p_x(x) p_y(y)}{p_x(x) p_y(y)} = 0
  \label{eq:E10}
\end{equation}

\end{document}